\section*{Tóm tắt}
\addcontentsline{toc}{section}{Tóm tắt}

Báo cáo trình bày bốn phần bắt buộc của Đề tài 9 \cite{de}: cài đặt từ đầu các
thành phần Transformer và hai hàm mục tiêu tinh chỉnh (Phần~A), chứng minh sáu
mệnh đề lý thuyết nền (Phần~B), bài tập tính toán nối lý thuyết với cài đặt
(Phần~C), và ứng dụng tinh chỉnh một mô hình ngôn ngữ thành trợ lý hỏi đáp thủ
tục hộ tịch -- chứng thực tại Hà Nội (Phần~D).

\textbf{Phần A \dalam.} Nhóm cài đặt bằng PyTorch, không dùng lớp attention hay
vòng huấn luyện dựng sẵn: chú ý tích vô hướng có chia tỷ lệ, chú ý nhiều đầu,
mặt nạ nhân quả, mã hóa vị trí hình sin, một khối decoder pre-LN, lớp LoRA có
hợp nhất trọng số, và hàm mất mát DPO. Mỗi cài đặt được đối chiếu với thư viện
chuẩn trên cùng đầu vào. Mô hình ngôn ngữ ký tự huấn luyện trên bản
\emph{Truyện Kiều} Trương Vĩnh Ký 1911 đạt perplexity validation
$\mathbf{10{,}1948 \pm 0{,}0628}$ (trung bình $\pm$ độ lệch chuẩn mẫu của 5 seed
độc lập).

\textbf{Phần B \dalam.} Sáu chứng minh đầy đủ giả thiết: phương sai tích vô hướng
bằng $d_k$ và vai trò hệ số $1/\sqrt{d_k}$; tính đẳng biến hoán vị của
self-attention cùng phản ví dụ cho thấy mặt nạ nhân quả phá tính chất này; biểu
diễn độ dịch vị trí bằng ma trận quay khối, kèm chứng minh mệnh đề \emph{thất bại}
ở số chiều lẻ; bốn tiểu mệnh đề LoRA; nghiệm duy nhất của bài toán reward có phạt
KL và dẫn xuất mất mát DPO qua Bradley--Terry; tương đương giữa cực tiểu
cross-entropy và cực tiểu $\KL(\pdata\|p_\theta)$ cùng quan hệ với perplexity.

\textbf{Phần C \dalam.} Bốn bài tập được giải và đối chiếu số với mã Phần~A: một
đầu attention nhân quả trên chuỗi 3 token; đếm $\mathbf{33{.}472}$ tham số và
$\mathbf{6{.}291{.}456}$ phép nhân--cộng cho một khối decoder theo cấu hình A.2 ở
độ dài 128, kèm phân tích hạng tử $O(n^2)$; số tham số LoRA ở Q/V của
Qwen2.5-1.5B-Instruct là $\mathbf{544{.}768}$, $\mathbf{1{.}089{.}536}$,
$\mathbf{2{.}179{.}072}$ cho $r=4,8,16$ cùng ước lượng bộ nhớ khi kết hợp QLoRA
4-bit; tính mất mát và gradient DPO, phân tích ảnh hưởng của $\beta$.

\textbf{Phần D \motphan.} Nhóm đã rà năm họ thủ tục từ nguồn chính văn, dựng quy
trình sinh QA bán tự động có cổng phát hành, và \textbf{chạy xong một lượt SFT
pilot thật} trên Kaggle: lô 24 QA đã duyệt 24/24, QLoRA $r=8$ ở Q/V, validation
PPL giảm $4{,}53\to3{,}46$ sau 2 epoch, peak allocated $2{,}402$ GiB, resume tái
lập đúng chỉ số. Tập chính đã lên \textbf{400 QA ứng viên} nhưng chưa phát hành.
Hỏi thử adapter pilot cho thấy hai lỗi thật: sai thời hạn, và \emph{không từ chối}
câu ngoài phạm vi. Chưa thu cặp so sánh, chưa huấn luyện DPO, chưa đánh giá ba mốc.
Mục~\ref{sec:phan-d} ghi rõ từng ngưỡng và khoảng cách còn lại.

\textbf{Kết luận trung thực.} Ba phần lý thuyết và cài đặt (A, B, C) có bằng chứng
kiểm chứng được trong kho mã. Phần ứng dụng còn thiếu các hạng mục phụ thuộc vào
con người thật: người khảo sát có văn bản đồng ý cho tập test, hai người chấm độc
lập cho tối thiểu 500 cặp ưu tiên, và người đánh giá mù cho so sánh cuối. Báo cáo
không trình bày số liệu giả định như kết quả đo được.
